\documentclass[11pt,reqno]{amsart}
\usepackage[T1]{fontenc}
\usepackage{lmodern}
\usepackage{microtype}
\usepackage{amsmath,amssymb,mathtools}
\usepackage[colorlinks=true,linkcolor=blue,citecolor=blue,urlcolor=blue]{hyperref}

\newtheorem{theorem}{Theorem}[section]
\newtheorem{lemma}[theorem]{Lemma}
\newtheorem{proposition}[theorem]{Proposition}
\newtheorem{corollary}[theorem]{Corollary}
\theoremstyle{definition}
\newtheorem{definition}[theorem]{Definition}
\newtheorem{remark}[theorem]{Remark}

\newcommand{\lcm}{\operatorname{lcm}}
\newcommand{\cR}{\mathcal R}
\newcommand{\cB}{\mathcal B}
\newcommand{\cE}{\mathcal E}
\newcommand{\cP}{\mathcal P}
\newcommand{\cQ}{\mathcal Q}
\newcommand{\Lc}{\mathcal L}
\newcommand{\E}{\mathbb E}
\newcommand{\PP}{\mathbb P}
\newcommand{\supp}{\operatorname{supp}}
\newcommand{\DT}{\operatorname{DT}}
\newcommand{\energy}{\operatorname{En}}
\newcommand{\Ho}{\mathcal H}
\newcommand{\cD}{\mathcal D}
\newcommand{\cM}{\mathcal M}
\newcommand{\cC}{\mathcal C}
\raggedbottom

% TODO(submission): replace "Anonymous" by the author metadata (kept anonymous for internal circulation).
\title[A sub-exponential Erd\H{o}s--Straus witness bound]{Large Erd\H{o}s--Straus witnesses II:\\
the multiplier modulus is infinitely often $\exp\big((\log p)^{1/5-o(1)}\big)$}
\author{Anonymous}
\date{Draft v3, 5 October 2026}

\begin{document}

\begin{abstract}
For a prime $p$ let $W(p)$ be the least modulus $M\equiv3\pmod4$ of a
``multiplier'' congruence that yields a solution of $4/p=1/x+1/y+1/z$.  In
earlier work we showed that $W(p)$ exceeds every fixed power of $\log p$ for
infinitely many primes $p$ in Mordell's hard classes.  Here we prove
$W(p)\ge\exp\big(c(\log p)^{1/5}(\log\log p)^{-1/5}\big)$ for infinitely many
hard primes, assuming Gallagher's prime number theorem for the family of all
characters and four results of Elsholtz and Tao on divisor sums of
polynomials; more precisely, the least hard prime with $W(p)>T$ is at most
$\exp(O((\log T)^5\log\log T))$.  Without the Elsholtz--Tao input we get
$\log W(p)\ge(1/\log2-o(1))\log\log p\cdot\log\log\log p$ infinitely often.
The class of one modulo a \emph{graded} set of prime powers leaves a residual
congruence system; the exponent at each prime is raised until the mass of the
system at that prime falls below a threshold proportional to the logarithm of
the prime power it would consume.  The local lemma then gives its avoiding
set positive Haar density, and the cost of the quarantine is controlled by a
moment of the number of prime-power divisors of the moduli, which we bound
with the Elsholtz--Tao divisor sums in progressions.  To transfer the density to primes we use
the one-sided $\ell^2$ sandwich of Bazzi and Razborov and a linear transfer
theorem based on Gallagher's theorem, as in the previous version.  The new
ingredient is a sharp energy bound for avoidance indicators: for any family
of cylinder events on a finite product probability space and any weights
$\lambda_v\ge1$ with $\prod_{v\in\supp E}\lambda_v\le2$ for every event $E$,
the Efron--Stein energies satisfy $\sum_U\prod_{v\in U}\lambda_v\|F^{=U}\|_2^2\le1$.
It replaces the binary encoding, Håstad's switching lemma and the
Linial--Mansour--Nisan argument of the previous version, and, applied to
the base-$\ell$ digits of the residues, it measures the cost of the
approximation in the size of the moduli instead of the number of primes.
\end{abstract}

\maketitle
% Work around a pdfTeX/microtype font-expansion warning (as in the first note).
\setbox0\hbox{\footnotesize\textsuperscript{1}}

\section{Introduction}

The Erd\H{o}s--Straus conjecture asserts that $4/n=1/x+1/y+1/z$ is solvable
in positive integers for every $n\ge2$.  It suffices to treat primes.  It has
been verified for $n\le10^{18}$ (\cite{MD}; \cite{Salez} for $10^{17}$), and
Mordell's identities settle every prime outside the six classes
$p\equiv1^2,11^2,13^2,17^2,19^2,23^2\pmod{840}$ \cite{Mordell} (cf.\
\cite{Oblath}).  We call the primes in these classes \emph{hard}.  For the
exceptional set see Vaughan \cite{Vaughan} and \cite{PW}; for the average
number of solutions see \cite{ET}.

For $M\equiv3\pmod4$ put $A_M=(M+1)/4$.  If $uvw=A_M$ and
$n\equiv-uv^{-1}\pmod M$, then $s=(nv+u)/M$ is a positive integer and
\begin{equation}\label{eq:ident}
 \frac4n=\frac1{suw}+\frac1{nsvw}+\frac1{nuvw}
\end{equation}
(direct verification; this is the classical multiplier family, cf.\
\cite{Vaughan,ET}).  Define
\begin{equation}\label{eq:W}
 W(n)=\min\{M\equiv3\ (4):\ n\equiv-uv^{-1}\ (\mathrm{mod}\ M)\ \text{for some}\ uvw=A_M\},
\end{equation}
with $W(n)=+\infty$ if no such $M$ exists.  Thus $W(p)\le T$ provides an
explicit representation with modulus at most $T$.  Every prime
$p\equiv1\pmod{\lcm\{M\le T\}}$ has $W(p)>T$ (Lemma~\ref{lem:one}), so
Linnik's theorem gives $W(p)\gg\log p$ infinitely often.  In
\cite{Paper1} we proved that for every fixed $k$,
$W(p)\ge(\log p)^{k}e^{-C_k\log\log p/\log\log\log p}$ for infinitely many
hard primes.  Here we prove a sub-exponential rate.

\begin{theorem}[Proved modulo Theorems~\ref{thm:G} and~\ref{thm:ETall}]\label{thm:main}
There is an absolute constant $c>0$ such that for infinitely many hard primes $p$
\[
 W(p)\ \ge\ \exp\Big(c\,\frac{(\log p)^{1/5}}{(\log\log p)^{1/5}}\Big).
\]
More precisely, for every sufficiently large $T$ there is a prime
$p\equiv1\pmod{840}$ with $p>T$, $W(p)>T$ and $\log p\le C(\log T)^{5}\log\log T$.
\end{theorem}

\begin{theorem}[Proved modulo Theorem~\ref{thm:G}]\label{thm:uncond}
For infinitely many hard primes $p$,
\[
 \log W(p)\ \ge\ \Big(\frac1{\log2}-o(1)\Big)\log\log p\cdot\log\log\log p .
\]
\end{theorem}

Theorem~\ref{thm:G} is Gallagher's prime number theorem for the family of
all primitive characters of conductor at most $Q_G$ \cite{Gallagher}, quoted
from \cite[Thm.~28.19]{MV3} together with the standard facts on exceptional
zeros (Theorem~\ref{thm:xz}).  No hypothesis on Siegel zeros is needed.
Theorem~\ref{thm:ETall} collects four results of Elsholtz and Tao
\cite{ET}: their Proposition~1.4, Theorem~7.1, Corollary~7.4 and the
estimate (7.10).  They are used only in \S\ref{sec:moment}, to bound the
total mass of the residual system by $O((\log T)^4)$ and a moment of the
number of prime-power divisors of its moduli by $O((\log T)^4\log\log T)$.
Without them the same argument gives Theorem~\ref{thm:uncond}.  All
constants are effective, except that we have not checked whether the
implied constants in \cite{ET} are.  We make no numerical claim.  We also
obtain a Haar-side bound (Corollary~\ref{cor:haar}): the set of integers
$n$ with $W(n)>T$ has density at least $\exp(-O((\log T)^5\log\log T))$,
modulo Theorem~\ref{thm:ETall}.

\emph{Versions.}  The first version of this note had exponent $1/14$; it used
the Thorner--Zaman theorem \cite{TZ} through the transfer theorem of
\cite{Paper1}.  The second had exponent $1/7$: it replaced Thorner--Zaman by
the linear transfer theorem (Theorem~\ref{thm:transfer} here, in a simpler
form), and bounded the approximation error by Håstad's switching lemma.  The
present version keeps the transfer theorem and the sandwich, and changes
three things: the switching lemma is replaced by an energy bound
(\S\ref{sec:energy}), the quarantine is graded (\S\ref{sec:system}), and the
cost of the quarantine is bounded through a moment estimate
(\S\ref{sec:moment}).  Håstad's switching lemma is no longer used.

\subsection*{The method}
Fix $T$.  We impose $n\equiv1$ modulo $Q=8\prod_\ell\ell^{a_\ell}$ for suitable
exponents $a_\ell\ge0$.  Then $W(n)>T$ whenever $n$ avoids a finite family
of congruence \emph{events} $n\equiv-4D\pmod{r}$, one for each surviving pair
$(M,D)$, where $r$ is the part of $M$ not absorbed by $Q$.  Unlike in
\cite{Paper1} and the previous version, a prime $\ell$ is not quarantined
all at once to $n\equiv1\pmod{\ell^{e}}$ with $\ell^e\le T<\ell^{e+1}$.  We
raise $a_\ell$ one step at a time, as long as the Haar mass $w_\ell$ of the
events that still involve $\ell$ exceeds $\theta_\ell=c(a_\ell+1)\log\ell/\log T$
(Lemma~\ref{lem:gradedq}).  Since every event comes from a modulus $M\le T$,
the total mass seen by one event is then at most $c$, which is what the local
lemma needs.  The cost $\log Q$ is at most $(\log T)/c$ times the moment
$\sum_{(M,D)}(g/M)\,h(M)$, where $g=\gcd(M,4D+1)$ and $h(M)=\sum_{\ell^i\mid M}1/i$.
We bound this moment by $O((\log T)^4\log\log T)$ (Theorem~\ref{thm:omega}), so
$\log Q\ll(\log T)^5\log\log T$.

To find a \emph{prime} in the avoiding set we need a function
$B=\sum_ic_i\mathbf 1[n\equiv b_i\ (d_i)]$ with $B\le\mathbf 1[\text{avoid}]$
pointwise, mean $\gg\delta$ (the Haar density of the avoiding set), Haar
$\ell^1$ norm $\E|B|$ comparable to its mean, and small moduli $d_i$.  The
linear transfer theorem (Theorem~\ref{thm:transfer}) then gives a prime $p$
with $\log p\ll\log Z$, where $Z$ is $Q$ times the largest $d_i$.  It is
Linnik's theorem with weights: the function $B(n)\mathbf 1[n\equiv1\ (Q)]$ is
expanded in characters, every coefficient is bounded by the Haar $\ell^1$
norm of $B$, whatever the number of cells, and Gallagher's theorem bounds
the contribution of all characters of conductor at most a power of $Z$ at once.

As in the previous version, $B$ is the one-sided $\ell^2$ sandwich of Bazzi
\cite{Bazzi}, in the form of Razborov \cite{Razborov} (Lemma~\ref{lem:brw}).
For events $E_1,\dots,E_m$ with indicators $A_i$ and arbitrary functions $u_j$,
\[
 B=1-\sum_iA_i\Big(1-\sum_{j<i}A_ju_j\Big)^2\ \le\ \mathbf 1[\text{no }E_i\text{ occurs}],
\]
and $\E[F-B]\le m^2\sum_j\PP(E_j)\,\E(F^{(j)}-u_j)^2$, where $F^{(j)}$ is the
indicator that no earlier event occurs, conditioned on $E_j$.  We need the
$u_j$ to be combinations of classes with small moduli.  The natural choice
is a truncation of the Efron--Stein decomposition of $F^{(j)}$, and the error
is then an Efron--Stein tail.  The previous version bounded this tail by
encoding residues in binary and applying Håstad's switching lemma
\cite{Hastad,LMN}.  Here we prove (Theorem~\ref{thm:C1}): if $F$ is the
avoidance indicator of any family of cylinder events on a finite product
probability space, and $\lambda_v\ge1$ are weights with
$\prod_{v\in\supp E}\lambda_v\le2$ for every event $E$, then
\[
 \sum_{U}\Big(\prod_{v\in U}\lambda_v\Big)\,\|F^{=U}\|_2^2\ \le\ 1 .
\]
With $\lambda_v=2^{1/k}$ this gives the tail $2^{-(t+1)/k}$ beyond level $t$ for
events on at most $k$ coordinates (Corollary~\ref{cor:tail}), with no
dependence on the alphabet sizes.  Applied to the base-$\ell$ digits of the
residues, with weights adapted to the moduli (Theorem~\ref{thm:filtration}),
it shows that truncating $F^{(j)}$ to moduli at most $e^{\tau}$ costs
$2^{-\tau/(2\log T)}$.  So moduli of size $\log d_i\ll\log T\cdot(S+\log T)$
suffice, where $S\ll(\log T)^4$ is the total mass.  Altogether
$\log Z\ll(\log T)^5\log\log T$.

\subsection*{Relation to the literature}
\emph{The energy bound.}  The nearest prior art known to us is Lecomte and
Tan \cite{LT}, who bound the Fourier coefficients of a DNF on $\{\pm1\}^n$ by
the probability that a set $S$ is covered by the terms satisfied at a random
point; this is the device of our Lemma~\ref{lem:cover}.  Their cover counts
are unsigned, and their degree concentration still comes from Håstad's
switching lemma, with an unspecified constant.  In Lemma~\ref{lem:cover} the
cover counts are signed, so that they vanish when a satisfied event misses
the set; the polarization identity (Lemma~\ref{lem:polar}) and the
deletion--contraction bound (Lemma~\ref{lem:theta}) then give the constant
$1$, on any finite product probability space.  In Boolean language
(Remark~\ref{rem:boolean}): every width-$k$ DNF $g:\{\pm1\}^n\to\{\pm1\}$
has $\sum_{|S|>t}\hat g(S)^2\le4\cdot2^{-(t+1)/k}$, under any product measure,
while the switching-lemma route gives $2\cdot2^{-t/(20k)}$ \cite[\S4.4]{OD}.
The base $2^{-1/k}$ is sharp over general product spaces
(Remark~\ref{rem:sharp}).  Theorem~\ref{thm:C1} and its corollaries are new to
us, but our literature search was partial (we checked \cite{LT}, \cite{OD} and
some lecture notes; we did not search the literature on switching-lemma
variants, Fourier growth or hypercontractivity with $\rho>1$
systematically), and we make no priority claim.

\emph{The rest.}  The sandwich (Lemma~\ref{lem:brw}) is Razborov's identity.
To our knowledge tools from the theory of bounded independence
(``polylogarithmic independence fools DNF and $\mathrm{AC}^0$''
\cite{Bazzi,Razborov,Braverman}) have not been used before to construct
sieve minorants for primes.  Gallagher's theorem is classically the engine of
Linnik's theorem on the least prime in a single progression.  Here it serves
as a transfer device for an arbitrary signed combination of progressions; the
cost is governed by the Haar $\ell^1$ norm of the minorant and its largest
modulus, not by the number of cells.  The nearest classical analogues are
bounds for the least prime in a union of residue classes, or in a Chebotarev
class of an abelian extension (Lagarias--Montgomery--Odlyzko \cite{LMO};
Thorner--Zaman \cite{TZcheb}); these concern nonnegative indicators.  We
have not searched systematically for an earlier use of Gallagher's theorem
in this role.  Earlier uses of the switching lemma in number theory known to
us run in the opposite direction: they show that the M\"obius function is
uncorrelated with bounded-depth circuits \cite{Green,Bourgain}.  The Haar
side, a local-lemma lower bound for the density of the uncovered set of a
congruence system, has relatives in the work on covering systems:
Filaseta, Ford, Konyagin, Pomerance and Yu \cite{FFKPY}, who to our knowledge
already use the local lemma for this purpose, Hough \cite{Hough} and
Balister, Bollob\'as, Morris, Sahasrabudhe and Tiba \cite{BBMST}.  We do not
use those results.  The moment bound of \S\ref{sec:moment} uses the
parametrisation and the divisor-sum estimates of Elsholtz and Tao
\cite{ET}; the splitting of the prime powers at the scale of the parameter
$c$ (Lemma~\ref{lem:rough}) is ours.

\smallskip\noindent\textbf{Status conventions.}  Every theorem header
carries a label.  \emph{Proved} means proved in full here.  \emph{Proved
modulo X} means proved here assuming the cited statement X, whose statement
we checked in the source but whose proof we did not check.  \emph{Cited}
marks external statements; \emph{Classical} marks standard results for
which we include a proof or a precise textbook reference.  The cited
inputs of the main theorems are Gallagher's theorem (Theorem~\ref{thm:G},
used together with the classical Landau--Page facts of
Theorem~\ref{thm:xz}) and the Elsholtz--Tao estimates
(Theorem~\ref{thm:ETall}).  In the labels, ``modulo Theorem~\ref{thm:G}''
includes Theorem~\ref{thm:xz}.  Classical results whose proof we include,
or which are standard textbook facts (the local lemma, the Efron--Stein
decomposition, Mertens' theorem, the divisor bound, the Rosser--Schoenfeld
bound for $\vartheta(x)$), are not listed in the labels.  \emph{Assessment}
marks heuristic comparisons, which are not claims.  Nothing here bears on
the truth of the Erd\H{o}s--Straus conjecture: a prime with large $W(p)$ may
well have other representations.

\smallskip\noindent\textbf{Notation.}  $T$ is large and $\Lc=\log T$;
$\log_2$ always denotes the binary logarithm (iterated logarithms are
written out as $\log\log$).  $\tau$ is the divisor function and
$\tau^*(x)=\max_{n\le x}\tau(n)$.  The letter $q$ denotes a prime power
$\ell^i$ with $i\ge1$, and $\sum_q$ is a sum over prime powers.  For
$\lambda=(\lambda_v)$ and a finite set $U$ we write $\lambda^U=\prod_{v\in U}\lambda_v$.

